\section{Appendix: Scope and Source Note}
\runninghead{Scope and Source Note}

This concise study treats the First Sunday of Advent (\latin{Dominica prima Adventus}, Sunday of the first class, \latin{Statio ad S.~Mariam maiorem}) in the 1962 typical \work{Missale Romanum} (\work{Proprium de tempore}, pp.~1--2; Trinity Preface, n.~1082, p.~293), under the 1960 \work{Codex rubricarum} (\work{Acta Apostolicae Sedis} [AAS]~52, 1960, pp.~597--729; \latin{Rubricae generales} [RG]~11, 15, 20, 91, 109, 111~a; \latin{Rubricae generales Missalis romani} [RGMR]~299, 427--429, 431--432, 434, 454, 457~d, 475--476, 480--482, 494~b, 498, 505, 507, 509--510; \work{Variationes}~II.13, IV.18). On Sunday, 29~November~2026 (General Roman Calendar), St.~Saturninus is omitted without seasonal extra orations. Local calendars govern particular celebrations; canonical authorization for 1962 use lies outside this scope. Postconciliar calendars, readings, prayers, and cycles are excluded.

All ten Latin proper texts and the Trinity Preface were visually collated against the CMAA 1962 facsimile PDF (collation control only) and corroborated verbatim in pre-1931 public-domain Missals (Pustet 1862, Tournai 1889, Tours Mame 1922, Cummiskey 1861) under 17~U.S.C.~\S103(b). Study English (never official 1962 liturgical English) follows public-domain Douay--Rheims (Challoner 1899 American ed., cross-checked against Gutenberg~1581, where Luke~21:29 reads ``a similitude'' vs.\ ``in a similitude'') and Lasance/Walsh, \work{The New Roman Missal} (1945, U.S.~public domain). Cited patristic, Doctoral, and liturgical editions are public domain except Wikisource Latin for St.~Gregory's \work{Homiliae in Evangelia} (CC~BY-SA~3.0, cited beside Migne PL~76, Bute~1908, and Oxford~1843). Each interpretation rests on at least two registered carrying authors (at least one Father of the Church); St.~Bernard of Clairvaux (canonized 1174, Doctor 1830) is cited as a supporting Doctor. Patristic exegesis remains distinct from editorial whole-Mass synthesis. For full texts, fuller exposition, chronology evidence notes, and search bounds, see the expansive study, \texttt{propers/verified.md}, and \texttt{research/scope.md}. This is a private study companion, not an official liturgical book, critical edition, or ecclesiastical approval.

\section{References}
\runninghead{References}

{\small
\begin{itemize}[itemsep=0pt,parsep=1pt,topsep=0.15em]
\item \work{Missale Romanum\dots{} editio typica} (Vatican City, 1962), \latin{Dominica prima Adventus}, pp.~1--2; \work{Praefationes}, n.~1082, p.~293; \work{Codex rubricarum} and \work{Variationes}~II.13, IV.18, \work{AAS}~52 (1960), pp.~597--729 (CMAA facsimile PDF, collation control only).
\item Pre-1931 public-domain Missals (17~U.S.C.~\S103(b)): Pustet (Ratisbon, 1862), pp.~1--3; Descl\'ee (Tournai, 1889); Mame (Tours, 1922), p.~2; E.~Cummiskey (Philadelphia, 1861; Trinity Preface).
\item Lasance, F.~X., and F.~A. Walsh, O.S.B., \work{The New Roman Missal in Latin and English} (New York: Benziger, rev.\ 1945), pp.~103, 105, 106, 773--774 (U.S.~public domain).
\item \work{Biblia Sacra iuxta Vulgatam Clementinam} (Ps.~24:1--4; Ps.~47:10; Ps.~79:3; Ps.~84:8, 13; Rom.~13:11--14; Luke~21:25--33) and Douay--Rheims (Challoner 1899 American ed., cross-checked against Gutenberg~1581; Ps.~24:1--4; Ps.~84:8, 13; Rom.~13:11--14; Luke~21:25--33; public domain).
\item Ambrose of Milan, St., \work{Expositio Evangelii secundum Lucam}, X.37--45 (PL~15, cols.~1812--1814; public domain).
\item Augustine of Hippo, St., \work{Enarrationes in Psalmos} on Ps.~24 (\S\S1--4) and Ps.~84 (\S\S9, 13--16; Coxe Ps.~85, \S\S6, 10--13, NPNF~I.8, 1888); \work{Sermo}~185.1 (\latin{In Natali Domini}~II, PL~38, cols.~997--998); \work{Confessions}~VIII.12.29 (NPNF~I.1); \work{De civitate Dei}~XX.14, 16, 24 (trans.\ M.~Dods, NPNF~I.2; public domain).
\item Bede, St., \work{In Lucae Evangelium Expositio}, VI, on Luke~21:28--29, 32--33 (PL~92; public domain).
\item Bernard of Clairvaux, St., \work{Sermones in Adventu Domini}, 2 (p.~19), 3.3--4 (pp.~25--26; PL~183, col.~45), 4.1 (pp.~32--33; PL~183, col.~47), 5.1--3 (pp.~40--42; PL~183, cols.~50--51), and \work{Missus est}, Hom.~1 (pp.~56, 58--60), in \work{Sermons for the Seasons}, vol.~1 (Dublin: Browne \& Nolan, 1921; public domain).
\item Cornelius a Lapide, \work{Commentaria in omnes Divi Pauli Epistolas} (Antwerp, 1614), pp.~195--198 (esp.~pp.~196--197) on Rom.~13:11--14 (public domain).
\item Cyril of Alexandria, St., \work{Commentary on S.~Luke}, Sermon~CXXXIX on Luke~21:25--28, trans.\ R.~Payne Smith (Oxford, 1859), pp.~653--654 (public domain).
\item Durandus, William, \work{Rationale divinorum officiorum}, VI.2--3 (Lyon, 1612, fols.~255v--261v; public domain).
\item Gregory the Great, St., \work{Homiliae in Evangelia}, Hom.~1, \S\S1--6 (PL~76, cols.~1077--1082; Wikisource Latin, CC~BY-SA~3.0; Bute, \work{Roman Breviary}, 1908; \work{Catena Aurea}, Oxford, 1843).
\item Gu\'eranger, Dom Prosper, \work{The Liturgical Year}, vol.~1 (\work{Advent}), trans.\ L.~Shepherd (Dublin: Duffy, 2nd~ed., 1870), pp.~127--138 (Communion, p.~137; public domain).
\item John Beleth, \work{Summa de ecclesiasticis officiis} LV (PL~202, col.~62) and LXIV (PL~202, cols.~71--72; public domain).
\item John Chrysostom, St., \work{Homilies on Romans}, Hom.~24 on Rom.~13:11--14 (NPNF~I.11; public domain).
\item Leo the Great, St., \work{Sermon}~12.4, \work{Sermon}~17.1, and \work{Sermon}~19.1 (NPNF~II.12, 1895; public domain).
\item Robert Bellarmine, St., \work{Explanatio in Psalmos}, trans.\ J.~O'Sullivan (Dublin, 1866), on Ps.~24:1--5, Ps.~79:2--4, and Ps.~84:8, 13 (public domain).
\item Rupert of Deutz, \work{De divinis officiis} III.2 (PL~170, col.~58; public domain).
\item Schuster, Bl.~Ildefonso, \work{The Sacramentary}, vol.~1, trans.\ A.~Levelis-Marke (London, 1924), pp.~321--323 (public domain).
\item Sicard of Cremona, \work{Mitrale} V.2 (PL~213, col.~205; public domain).
\item Thomas Aquinas, St., \work{Postilla super Psalmos} on Ps.~24, \work{Opera omnia}, vol.~14 (Parma, 1863), pp.~229--230; \work{Catena Aurea in Lucam} on Luke~21:28 (Oxford, 1843/1865; public domain).
\end{itemize}
}
